\chapter{Calculer et justifier les dérivées}
\section*{Exercice 9.1 -- Somme et puissances}
Le domaine est $\R\setminus\{0\}$. Les règles de somme et des puissances donnent $f'(x)=12x^3+4x^{-3}$. La constante $7$ a une dérivée nulle.

\section*{Exercice 9.2 -- Produit}
Le domaine de dérivabilité est $]-2,+\infty[$. La règle du produit donne $2x\sqrt{x+2}+(x^2+1)/(2\sqrt{x+2})=(5x^2+8x+1)/(2\sqrt{x+2})$. À $-2$, la formule n'est pas définie et ne peut être prolongée par simple substitution.

\section*{Exercice 9.3 -- Quotient}
Pour $x\ne3$, $f'(x)=[2(x-3)-(2x+1)]/(x-3)^2=-7/(x-3)^2$. La réécriture $f(x)=2+7/(x-3)$ donne immédiatement la même dérivée. Deux expressions peuvent réduire le risque d'erreur.

\section*{Exercice 9.4 -- Composition affine}
La dérivée de la fonction extérieure est multipliée par celle de l'expression intérieure : $g'(x)=15(3x-2)^4$ sur $\R$. Pour $x\ne-1/2$, $h'(x)=-2/(2x+1)^2$. Oublier les facteurs $3$ ou $2$ serait oublier la composition.

\section*{Exercice 9.5 -- Démontrer une règle}
Ajouter et retrancher $f(a+h)g(a)$ donne $f(a+h)[g(a+h)-g(a)]+g(a)[f(a+h)-f(a)]$. Diviser par $h$, puis passer à la limite. La dérivabilité de $f$ implique sa continuité, donc $f(a+h)\to f(a)$. On obtient $(fg)'(a)=f(a)g'(a)+g(a)f'(a)$.

\section*{Exercice 9.6 -- Une erreur à réfuter}
Prendre $f(x)=x^2$, $g(x)=x$ sur $]0,+\infty[$. Comme $f/g=x$, sa dérivée vaut $1$, alors que $f'/g'=2x$. En $x=1$, on obtient $1\ne2$. La vraie règle contient $f'g-fg'$ au numérateur et $g^2$ au dénominateur.
